\documentclass[12pt, twoside]{article}
\input{../preamble}

\fancyhead[LE]{\thepage}
\fancyhead[RO]{Markscheme}
\fancyhead[LO]{La Scuola d'Italia / Dr.\ Huson / IB Math AA SL \\* 12 May 2026}
\fancyfoot[C]{\thepage}

\pgfplotsset{compat=1.18}

\begin{document}
\subsubsection*{11.17 Classwork: Cumulative Frequency and Normal Models --- Markscheme}

\begin{enumerate}[itemsep=0.3cm]

%--- Problem 1: AA SL 2023 Nov P1 TZ2 Q7 [17 marks] ---
\item[1.] \makebox[\linewidth][s]{[Maximum mark: 17]\hfill {\small AA SL 2023 Nov P1 TZ2 Q7}}

\textbf{(a)(i)} Find $p$. \hfill [1]

$p=60-3-18-30=9$ \hfill \textbf{A1}

\medskip

\textbf{(a)(ii)} Write down the modal class. \hfill [1]

$600<n\leq 800$ \hfill \textbf{A1}

\emph{Do not accept $700$.}

\medskip

\textbf{(b)(i)} Estimate the median number of tickets sold. \hfill [1]

The median is the 30th performance. From the cumulative frequency graph,
the median is $600$ tickets. \hfill \textbf{A1}

\medskip

\textbf{(b)(ii)} Estimate the number of performances where at least
$80\%$ of the tickets were sold. \hfill [3]

$80\%$ of $800$ is $640$. \hfill \textbf{A1}

From the cumulative frequency graph, about $40$ performances sold fewer
than $640$ tickets. \hfill \textbf{A1}

Therefore $60-40=20$ performances sold at least $80\%$ of the tickets.
\hfill \textbf{A1}

\medskip

\textbf{(c)(i)} State one disadvantage of surveying only the first
$5\%$ of the audience as they leave. \hfill [1]

Any reasonable answer identifying bias, with a reason. For example:
the first people leaving may come from the same section, group, or
priority seating area, so the sample may not represent the whole
audience. \hfill \textbf{A1}

\medskip

\textbf{(c)(ii)} Describe how to collect feedback from $5\%$ of the
audience using systematic sampling. \hfill [2]

Select people at regular intervals. \hfill \textbf{A1}

Since $5\%=\dfrac{1}{20}$, survey every 20th person. \hfill \textbf{A1}

\medskip

\textbf{(c)(iii)} State the sampling method which should be used if the
survey is to be representative of the number of children and adults.
\hfill [1]

Quota sampling \hfill \textbf{A1}

\medskip

\textbf{(d)(i)} Estimate the number of people who spent between $\$3$
and $\$25$. \hfill [2]

From the box plot, $\$3$ is the lower quartile and $\$25$ is the maximum,
so this interval contains $75\%$ of the audience. \hfill \textbf{M1}

$0.75\times 36\,000=27\,000$ people. \hfill \textbf{A1}

\medskip

\textbf{(d)(ii)} Half the audience spent less than $\$a$. Estimate $a$.
\hfill [1]

The median amount spent is $\$7$, so $a=7$. \hfill \textbf{A1}

\medskip

\textbf{(e)(i)} Calculate the mean number of tickets expected this year.
\hfill [3]

Last year the total number of tickets sold was $36\,000$, so the old
mean was
\[
\frac{36\,000}{60}=600.
\]
Old mean is $600$ tickets. \hfill \textbf{A1}

Adding $17$ tickets to each performance adds $17$ to the mean.
\hfill \textbf{M1}

New mean $=600+17=617$ tickets. \hfill \textbf{A1}

\emph{Alternative: compute
$\dfrac{36\,000+60(17)}{60}=617$.}

\medskip

\textbf{(e)(ii)} State the effect on the variance. \hfill [1]

Adding the same constant to every value has no effect on the variance.
\hfill \textbf{A1}

\newpage

%--- Problem 2: AA SL 2023 Nov P2 TZ1 Q9 [16 marks] ---
\item[2.] \makebox[\linewidth][s]{[Maximum mark: 16]\hfill {\small AA SL 2023 Nov P2 TZ1 Q9}}

\textbf{(a)} Find $P(82.4 < H < 87.3)$. \hfill [2]

Recognize that the three probabilities sum to $1$. \hfill \textbf{M1}

\[
0.213+P(82.4<H<87.3)+0.409=1
\]

\[
P(82.4<H<87.3)=0.378
\]
\hfill \textbf{A1}

\medskip

\textbf{(b)} Find $\mu$ and $\sigma$. \hfill [5]

Use inverse normal with $P(H<82.4)=0.213$ and
$P(H<87.3)=1-0.409=0.591$. \hfill \textbf{M1}

\[
\mu+z_{1}\sigma=82.4,\qquad \mu+z_{2}\sigma=87.3
\]
where $z_1=\text{invNorm}(0.213)=-0.796055\ldots$ and
$z_2=\text{invNorm}(0.591)=0.230118\ldots$. \hfill \textbf{A1A1}

Attempt to solve the two equations in $\mu$ and $\sigma$. \hfill \textbf{M1}

\[
\mu=86.2011\ldots,\qquad \sigma=4.77502\ldots
\]

\[
\mu=86.2,\qquad \sigma=4.78
\]
\hfill \textbf{A1}

\emph{Accept equivalent GDC setup using normal CDF or inverse normal.}

\medskip

\textbf{(c)(i)} Find the probability that exactly $32$ plants are ready
to harvest. \hfill [2]

Recognition of binomial distribution. \hfill \textbf{M1}

\[
X\sim B(100,0.409)
\]

\[
P(X=32)=0.0157931\ldots=0.0158
\]
\hfill \textbf{A1}

\medskip

\textbf{(c)(ii)} Given that fewer than $44$ plants are ready to harvest,
find the probability that exactly $32$ plants are ready to harvest.
\hfill [4]

\[
P(X<44)=0.702975\ldots
\]
\hfill \textbf{A1}

Recognition of conditional probability in context. \hfill \textbf{M1}

\[
P(X=32\mid X<44)=\frac{P(X=32)}{P(X<44)}
\]
\hfill \textbf{A1}

\[
\frac{0.0157931\ldots}{0.702975\ldots}=0.0224661\ldots=0.0225
\]
\hfill \textbf{A1}

\emph{FT: award marks for correct conditional setup using the
candidate's binomial probability from part (c)(i).}

\medskip

\textbf{(d)} Find $d$. \hfill [3]

Since the IQR is $4.52$, the quartiles are
\[
Q_1=92.8-2.26=90.54,\qquad Q_3=92.8+2.26=95.06.
\]
One correct quartile value may be seen. \hfill \textbf{A1}

Use the fact that $P(F<90.54)=0.25$ or $P(F<95.06)=0.75$.
\hfill \textbf{M1}

\[
\frac{90.54-92.8}{d}=-0.674489\ldots
\]

\[
d=3.35068\ldots=3.35
\]
\hfill \textbf{A1}

\emph{Accept equivalent GDC method solving for $d$.}

\end{enumerate}
\end{document}
